
\section{\BookRecordIndex{T0188}截断函数}\label{sec:10-4}

迄今为止，我们用任意光滑\BookRecordIndex{T0188}截断函数 $\rho$ 定义\BookRecordIndex{T0394}预粘合 $u^R$。其余构造也都将通过拼合
分别定义在区域 $\{|z|\ge1/R\}$ 和 $\{|z|\le1/R\}$ 上的对象来完成。然而，为使所需
估计成立，必须使用一个经过精细调节的\BookRecordIndex{T0188}截断函数，记为 $\beta$。本节中所有积分都理解为
关于 Lebesgue 测度的积分。读者初读时完全可以跳过本节，待读到命题~\ref{prop:10-5-1} 的证明中
$\beta$ 的各种性质如何被使用后再返回。

\begin{lemma}\label{lem:10-4-1}
对任意两个实数 $0<\delta<\varepsilon$，存在 Lipschitz 连续\BookRecordIndex{T0188}截断函数
$\beta=\beta_{\delta,\varepsilon}:\mathbb R^2\to[0,1]$，使得
\[
  \beta(z)=
  \begin{cases}
    1,& |z|\le\delta,\\
    0,& |z|\ge\varepsilon,
  \end{cases}
  \qquad
  \int_{\mathbb R^2}|\nabla\beta|^2
  =\frac{2\pi}{\log(\varepsilon/\delta)}.
\]
\end{lemma}

\begin{proof}
定义 $\beta:\mathbb R^2\to[0,1]$：当 $|z|\le\delta$ 时令 $\beta(z):=1$，当
$|z|\ge\varepsilon$ 时令 $\beta(z):=0$，而在中间区域令
\begin{equation}\label{eq:10-4-1}
  \beta(z):=
  \frac{\log\varepsilon-\log|z|}{\log\varepsilon-\log\delta},
  \qquad \delta\le|z|\le\varepsilon.
\end{equation}
于是 $|\nabla\beta(z)|=1/(|z|\log(\varepsilon/\delta))$，从而
\[
  \begin{aligned}
  \int_{\mathbb R^2}|\nabla\beta(z)|^2
  &=\int_{\delta\le|z|\le\varepsilon}
    \frac{1}{|z|^2\log(\varepsilon/\delta)^2}\\
  &=\int_\delta^\varepsilon
    \frac{2\pi}{s\log(\varepsilon/\delta)^2}\,ds
   =\frac{2\pi}{\log(\varepsilon/\delta)}.
  \end{aligned}
\]
引理得证。
\end{proof}

要强调的是，引理~\ref{lem:10-4-1} 中 $\nabla\beta$ 的 $L^2$ 范数只依赖于比值
$\varepsilon/\delta$。这说明在二维中，导数的 $L^2$ 范数具有共形不变性。此外，若将常数
$2\pi/\log(\varepsilon/\delta)$ 换成任意更大的常数，显然可以把\BookRecordIndex{T0188}截断函数 $\beta$ 选成
光滑的。因此，$\mathbb R^2$ 上存在一列紧支函数，它们在 $W^{1,2}$ 范数下收敛到零，
但不在 $L^\infty$ 范数下收敛到零（固定常数 $\varepsilon>0$，让 $\delta$ 趋于零即可）。
这是索伯列夫估计的临界情形。当 $p>2$ 时存在嵌入
$W^{1,p}(\mathbb R^2)\hookrightarrow C(\mathbb R^2)$，所以若把前一引理中
$\nabla\beta$ 的 $L^2$ 范数换成 $p>2$ 时的 $L^p$ 范数，其断言便不再成立。下面将看到，
在二维\BookRecordIndex{T0254}粘合构造中证明一阶算子满射时，所需的恰好是 $L^2$ 范数。$n$ 维中的类似论证要求
$\nabla\beta$ 的 $L^n$ 范数很小，这同样是索伯列夫临界情形。$n=4$ 的情形与反自对偶
瞬子的\BookRecordIndex{T0254}粘合有关；Donaldson--Kronheimer \srcite{88, pp.~287--295} 对此作了解释。

现在令 $\beta$ 由~\eqref{eq:10-4-1} 给出，并取
$\xi\in W^{1,q}(\mathbb R^2)$，其中 $1\le q<2$。于是，对某个合适的常数 $r\ge2$，有
如下估计\sourcefootnote{原书在式~\eqref{eq:10-4-2} 最后一项中误印为
$2\pi c/\log(\varepsilon/\delta)$。由引理~\ref{lem:10-4-1}，
$\|\nabla\beta\|_{L^2}=\sqrt{2\pi/\log(\varepsilon/\delta)}$，故该项唯一应为这里的
平方根形式。}：
\begin{equation}\label{eq:10-4-2}
  \begin{aligned}
  \|\nabla\beta\cdot\xi\|_{L^q}
  &\le\|\nabla\beta\|_{L^2}\|\xi\|_{L^r}
   \le c\|\nabla\beta\|_{L^2}\|\xi\|_{W^{1,q}}\\
  &\le c\sqrt{\frac{2\pi}{\log(\varepsilon/\delta)}}
  \|\xi\|_{W^{1,q}}.
  \end{aligned}
\end{equation}
这里，第一个不等式使用 Hölder 不等式
\[
  \left(\int|uv|^q\right)^{1/q}
  \le
  \left(\int|u|^s\right)^{1/s}
  \left(\int|v|^r\right)^{1/r},
  \qquad
  \frac1s+\frac1r=\frac1q,
\]
其中 $s=2$。因此必须取 $r:=2q/(2-q)$；由定理 B.1.12，有索伯列夫嵌入
$W^{1,q}\hookrightarrow L^r$，这给出上面的第二个不等式。第三个不等式由
引理~\ref{lem:10-4-1} 得出。

由~\eqref{eq:10-4-2} 可知，只要 $q<2$，乘法算子
\[
  W^{1,q}(\mathbb R^2)\longrightarrow L^q(\mathbb R^2),
  \qquad \xi\longmapsto\nabla\beta_{\delta,\varepsilon}\cdot\xi
\]
就在 $\delta\to0$ 时于算子范数下收敛到零。这个原样的断言对 $q>2$ 不再成立。不过，
在 $q>2$ 的情形，可以把这些算子限制到所有在原点消失的函数所成子空间。下一个引理说明，
当 $q>2$ 时，这些受限算子在 $\delta\to0$ 时仍于算子范数下收敛到零。记
\[
  B_r:=\{z\in\mathbb R^2\mid |z|\le r\},
  \qquad
  A(r,R):=\{z\in\mathbb R^2\mid r\le|z|\le R\}.
\]

\begin{lemma}\label{lem:10-4-2}
对每个 $p>2$，存在具有如下性质的常数 $c>0$。给定 $0<\delta<\varepsilon$，用
\eqref{eq:10-4-1} 定义 $\beta:\mathbb R^2\to[0,1]$。则对每个满足
$\xi(0)=0$ 的 $\xi\in W^{1,p}(B_\varepsilon)$，都有
\[
  \|\nabla\beta\cdot\xi\|_{L^p(A(\delta,\varepsilon))}
  \le
  \frac{c}{\log(\varepsilon/\delta)^{1-1/p}}
  \|\xi\|_{W^{1,p}(B_\varepsilon)}.
\]
\end{lemma}

\begin{proof}
由引理~\ref{lem:10-4-1} 的证明，$\nabla\beta$ 的支集包含在\BookRecordIndex{T0170}环域
$A(\delta,\varepsilon)$ 中，并且当 $\delta\le|z|\le\varepsilon$ 时，
\[
  |\nabla\beta(z)|=\frac{1}{|z|\log(\varepsilon/\delta)}.
\]
此外，由定理 B.1.11，每个 $\xi\in W^{1,p}(\mathbb R^2)$ 都以指数
$\mu=1-2/p$ Hölder 连续。更准确地说，存在只依赖于 $p>2$ 的常数 $c>0$，使得
\begin{equation}\label{eq:10-4-3}
  z_0,z_1\in B
  \quad\Longrightarrow\quad
  |\xi(z_1)-\xi(z_0)|
  \le c\|\xi\|_{W^{1,p}(B)}|z_1-z_0|^{1-2/p}
\end{equation}
对每个球 $B\subset\mathbb R^2$ 和每个 $\xi\in W^{1,p}(B)$ 都成立（见
引理 B.1.16 的证明）。因此，每个在 $z=0$ 处消失的
$\xi\in W^{1,p}(B_\varepsilon)$ 都满足估计
\[
  \begin{aligned}
  \int_{\delta\le|z|\le\varepsilon}
  |\nabla\beta(z)|^p|\xi(z)|^p
  &\le
  \frac{c^p}{\log(\varepsilon/\delta)^p}
  \int_{\delta\le|z|\le\varepsilon}\frac{1}{|z|^2}
  \|\xi\|_{W^{1,p}(B_\varepsilon)}^p\\
  &=\frac{2\pi c^p}{\log(\varepsilon/\delta)^p}
  \int_\delta^\varepsilon\frac1s\,ds\,
  \|\xi\|_{W^{1,p}(B_\varepsilon)}^p\\
  &=\frac{2\pi c^p}{\log(\varepsilon/\delta)^{p-1}}
  \|\xi\|_{W^{1,p}(B_\varepsilon)}^p.
  \end{aligned}
\]
引理得证。
\end{proof}

我们还需要任意光滑\BookRecordIndex{T0188}截断函数 $\rho:\mathbb C\to[0,1]$ 的如下结论；这里 $\rho$ 在
$|z|\le1$ 时等于零，在 $|z|\ge2$ 时等于一。

\begin{lemma}\label{lem:10-4-3}
设 $p>2$，且 $\xi\in W^{1,p}(B_1)$ 满足 $\xi(0)=0$。记
$\rho_R(z):=\rho(Rz)$。则
\[
  \lim_{R\to\infty}\|\xi-\rho_R\xi\|_{W^{1,p}(B_1)}=0.
\]
\end{lemma}

\begin{proof}
我们证明存在常数 $c>0$，使得
\begin{equation}\label{eq:10-4-4}
  \|\xi-\rho_R\xi\|_{W^{1,p}(B_1)}
  \le c\|\xi\|_{W^{1,p}(B_{2/R})}.
\end{equation}
首先注意，若把 $W^{1,p}$ 范数换成 $L^p$ 范数，则这个估计取 $c=1$ 即成立。其次，
\[
  \nabla(\xi-\rho_R\xi)
  =(1-\rho_R)\nabla\xi-(\nabla\rho_R)\xi.
\]
第一项的 $L^p$ 范数以 $c=1$ 满足所需不等式。为估计第二项，注意
$\|\nabla\rho_R\|_{L^\infty}=Rc_0$，其中
$c_0:=\|\nabla\rho\|_{L^\infty}$。此外，由~\eqref{eq:10-4-3}，存在常数 $c_1>0$，使得
\[
  |\xi(z)|
  \le c_1|z|^{1-2/p}\|\xi\|_{W^{1,p}(B_{2/R})},
  \qquad |z|\le2/R.
\]
因此
\[
  \begin{aligned}
  \|(\nabla\rho_R)\xi\|_{L^p}
  &\le c_0c_1R
  \left(\int_{|z|\le2/R}|z|^{p-2}\right)^{1/p}
  \|\xi\|_{W^{1,p}(B_{2/R})}\\
  &=c_0c_1\left(\frac{2^{p+1}\pi}{p}\right)^{1/p}
  \|\xi\|_{W^{1,p}(B_{2/R})}.
  \end{aligned}
\]
这证明了~\eqref{eq:10-4-4}，也证明了引理。
\end{proof}
